\documentclass{situationist-preprint}

\title{Agenda-Setting in Academic Publishing}
\msid{POL-10-b}
\msfamily{POL-10}
\msdiscipline{Political Philosophy}
\mstype{Companion manuscript}
\mskeywords{power; agenda-setting; peer review; academic journals; research funding; preference formation}

\begin{document}
\maketitle

\begin{abstract}
POL-10 concludes that power is the capacity to shape agendas and preferences as well as outcomes. We apply the conclusion to academic publishing. Referees and editors exercise power along the decision route when they accept or reject; editors and funders exercise it along the agenda route when they determine which work is considered at all and which questions are posed; and the institutions of training and reward exercise it along the preference route by shaping what researchers want to work on. By route-neutrality, the second and third are power as much as the first. Whether any of them is objectionable is a further question, about domination, which turns on whether the power is checked and answerable.
\end{abstract}

\section{Three routes in publishing}

The outcomes at stake in academic publishing are which claims enter the published record, with what prominence, and so which claims a field treats as established. POL-10-a distinguishes three routes by which an agent may have the capacity to make a difference to outcomes (Definition~2.2 of POL-10-a).

\emph{The decision route.} Referees and editors decide whether a submitted manuscript is published. This is the visible face of power in publishing, and the one on which discussion of peer review concentrates: its reliability, its consistency, and its biases. A classic study resubmitted published articles to the journals that had published them and found most rejected, few detected as resubmissions \citep{Peters1982}.

\emph{The agenda route.} Before any decision on the merits, editors determine which submissions are sent for review; journals determine by their scope which topics are within the range of submissions; special issues and invited contributions place questions on the agenda; and funding bodies, by the calls they issue, determine which research questions are pursued at all. A topic that falls outside every journal's scope and every funder's priorities is kept off the agenda as reliably as one rejected on the merits, which is the face of power that \citet{Bachrach1962} called non-decision-making.

\emph{The preference route.} Doctoral training, hiring and promotion criteria and citation-based rewards shape what researchers want to work on and which questions seem to them worth asking. Prominence breeds prominence, so that the work of those already recognized attracts disproportionate attention \citep{Merton1968}; and what a field takes as its exemplary problems shapes the problems its members find interesting \citep{Kuhn1962}.

\section{Route-neutrality applied}

\begin{proposition}\label{p:power}
Editors and funders who control which questions reach the stage of decision, and the institutions that shape researchers' preferences, have power over the published record, to the extent that their control over which claims enter it matches that of referees deciding on submissions.
\end{proposition}

\begin{proof}
By Principle~P2 of POL-10-a, capacities giving the same control over which outcomes come about are power alike, whatever their route. A funder that declines to fund a line of research secures its absence from the record as reliably as referees who would reject its results; an editor who desk-rejects a topic secures the same outcome as reviewers who would reject each paper on it. And a training regime that leaves no one wanting to pursue a question secures its absence through researchers' own choices. Each instantiates Assumption~2.3 or~2.4 of POL-10-a in this domain, so Theorem~3.2 of POL-10-a applies.
\end{proof}

Two consequences follow. First, the power of the agenda and preference routes is typically greater than that of the decision route, since it operates before any manuscript exists, and attention confined to the reliability of peer review misses most of the power in the system. Second, by Principle~P1 of POL-10-a, power is a capacity: a funder's power over a field does not depend on its ever refusing an application, since researchers write the applications they expect to succeed.

\section{Power and domination}

It does not follow that the power of editors, funders or training is objectionable. By Corollary~3.3 of POL-10-a, power does not require that it be exercised against anyone's real interests, and much agenda-setting in publishing serves researchers' interests: a field without editors would have no shared agenda. Whether a given exercise is objectionable is a question for the theory of domination, which asks whether the power is checked, answerable and justifiable to those subject to it \citep{Pettit1997}. The institutional remedies that this suggests, such as transparent scope statements, rotation of editorial and funding panels, and avenues of appeal, are remedies for unaccountable power, not for power as such.

\begin{objection}
Editors and funders merely respond to the community's judgement of what matters; the power is the community's, not theirs.
\end{objection}

\begin{reply}
That the community's preferences constrain editors does not show that editors lack power; it shows that power is distributed. And the community's preferences are in part shaped by the same institutions, which is the preference route. Distributed power is still power, and its accountability is still in question.
\end{reply}

\section{Conclusion}

Academic publishing illustrates the conclusion of POL-10. Power over the published record lies not only in the decisions of referees but in the setting of agendas by editors and funders and in the shaping of researchers' preferences. Whether that power is legitimate depends on whether it is answerable to those it affects.

\begin{thebibliography}{9}
\bibitem[Bachrach and Baratz(1962)]{Bachrach1962} Bachrach, P. and Baratz, M.~S. (1962). Two faces of power. \emph{American Political Science Review}, 56(4), 947--952.
\bibitem[Kuhn(1962)]{Kuhn1962} Kuhn, T.~S. (1962). \emph{The Structure of Scientific Revolutions}. University of Chicago Press.
\bibitem[Pettit(1997)]{Pettit1997} Pettit, P. (1997). \emph{Republicanism: A Theory of Freedom and Government}. Oxford University Press.
\bibitem[Merton(1968)]{Merton1968} Merton, R.~K. (1968). The Matthew effect in science. \emph{Science}, 159(3810), 56--63.
\bibitem[Peters and Ceci(1982)]{Peters1982} Peters, D.~P. and Ceci, S.~J. (1982). Peer-review practices of psychological journals: The fate of published articles, submitted again. \emph{Behavioral and Brain Sciences}, 5(2), 187--195.
\end{thebibliography}

\end{document}
